\documentclass[12pt]{article}
\usepackage[T2A]{fontenc}
\usepackage[utf8]{inputenc}
\usepackage[russian]{babel}
\usepackage{amsmath}
\usepackage{booktabs}
\usepackage{geometry}
\geometry{margin=2cm}

\begin{document}

\begin{center}
\textbf{Лабораторная работа №2}\\
\textbf{Изучение алгоритма RSA}\\
\textbf{М26-545 Губанов Алексей Дмитриевич Вариант 11}
\end{center}

\section{Цель работы}
Освоить механизм шифрования и дешифрования данных в криптографической системе с открытыми ключами RSA.

\section{Постановка задачи}
Разработать программу, осуществляющую генерацию ключей RSA по заданным \(p\) и \(q\), шифрование и дешифрование сообщения. Предусмотреть проверку взаимной простоты чисел, вывод ключей и зашифрованного текста. Оформить результаты в виде отчёта.

\section{Исходные данные}
По варианту 11:
\[
p=241,\qquad q=307.
\]
Оба числа являются простыми. Тогда
\[
n=pq=241\cdot 307=73987,
\]
\[
\varphi(n)=(p-1)(q-1)=240\cdot 306=73440.
\]
Минимальное \(e>1\), взаимно простое с \(\varphi(n)\), равно
\[
e=7,\qquad \gcd(7,73440)=1.
\]
Секретная экспонента:
\[
d=e^{-1}\bmod \varphi(n)=7^{-1}\bmod 73440=20983.
\]
Проверка:
\[
7\cdot 20983=146881=2\cdot 73440+1.
\]

\section{Ключи}
\begin{center}
\begin{tabular}{ll}
\toprule
Параметр & Значение \\
\midrule
\(p\) & 241 \\
\(q\) & 307 \\
\(n=pq\) & 73987 \\
\(\varphi(n)\) & 73440 \\
\(e\) & 7 \\
\(d\) & 20983 \\
Открытый ключ & \((e,n)=(7,73987)\) \\
Секретный ключ & \((d,n)=(20983,73987)\) \\
\bottomrule
\end{tabular}
\end{center}

\section{Алгоритм работы программы}
\begin{enumerate}
    \item Проверить простоту \(p\) и \(q\).
    \item Вычислить \(n=pq\) и \(\varphi(n)=(p-1)(q-1)\).
    \item Выбрать минимальное \(e>1\), для которого \(\gcd(e,\varphi(n))=1\).
    \item Найти \(d\) из условия \(ed\equiv 1\pmod{\varphi(n)}\).
    \item Закодировать сообщение в UTF-8 и разбить на блоки.
    \item Для каждого блока \(m\) вычислить \(c=m^e\bmod n\).
    \item Для дешифрования вычислить \(m=c^d\bmod n\).
    \item Восстановить исходную строку и проверить совпадение.
\end{enumerate}

Для \(n=73987\) выполнено \(256^2=65536<n\), поэтому используется размер блока 2 байта.

\section{Результаты работы программы}
В качестве примера использовано сообщение
\[
M=\text{«RSA»}.
\]
После UTF-8-кодирования и разбиения на 2-байтовые блоки получены числа:
\[
m_1=21075,\qquad m_2=16640.
\]
Шифрование:
\[
c_1=21075^7\bmod 73987=69939,
\]
\[
c_2=16640^7\bmod 73987=24011.
\]
Зашифрованный текст:
\[
C=[69939,\ 24011].
\]
После дешифрования:
\[
M'=\text{«RSA»}.
\]
Исходное и расшифрованное сообщения совпадают.

\begin{center}
\begin{tabular}{ll}
\toprule
Сообщение & RSA \\
Блоки \(m\) & 21075, 16640 \\
Шифртекст \(c\) & 69939, 24011 \\
Расшифровка & RSA \\
Проверка & успешно \\
\bottomrule
\end{tabular}
\end{center}

\section{Анализ результатов}
Ключи сгенерированы корректно. Минимальное допустимое \(e=7\) взаимно просто с \(\varphi(n)\). Секретная экспонента \(d=20983\) удовлетворяет условию \(ed\equiv1\pmod{\varphi(n)}\). Шифрование и дешифрование выполнены без потерь.

\section{Выводы}
В ходе работы реализован алгоритм RSA. Получены открытый и секретный ключи, выполнены шифрование и дешифрование сообщения. Корректность восстановления подтверждена. Стойкость RSA основана на сложности факторизации числа \(n=pq\).

\section{Ответы на контрольные вопросы}

\begin{enumerate}
    \item \textbf{Что такое однонаправленные функции?}  
    Это функции \(f(x)\), для которых при известном \(x\) значение \(f(x)\) вычисляется просто, а при известном \(y=f(x)\) найти \(x\) практически невозможно.

    \item \textbf{Основные свойства однонаправленных функций с потайным ходом.}  
    Такая функция является однонаправленной, но при наличии секретной дополнительной информации (потайного хода) обращение функции становится вычислительно простым.

    \item \textbf{Какие числа называются взаимно простыми?}  
    Числа \(a\) и \(b\) называются взаимно простыми, если \(\gcd(a,b)=1\).

    \item \textbf{Как реализуется программное возведение в степень для больших чисел?}  
    Используется алгоритм бинарного возведения в степень (square-and-multiply). На каждом шаге выполняется умножение и приведение по модулю \(n\), что даёт сложность \(O(\log e)\) модульных умножений.

    \item \textbf{На чём основана криптостойкость алгоритма RSA?}  
    На вычислительной сложности факторизации большого числа \(n=pq\). Если разложить \(n\) на \(p\) и \(q\), можно вычислить \(\varphi(n)\) и секретную экспоненту \(d\).

    \item \textbf{Каковы достоинства и недостатки асимметричных алгоритмов?}  
    Достоинства: не требуется предварительная передача секретного ключа, возможны цифровая подпись и аутентификация.  
    Недостатки: низкая скорость по сравнению с симметричными алгоритмами, большие размеры ключей, необходимость инфраструктуры открытых ключей.

    \item \textbf{В чём заключается основная проблема симметричных криптосистем?}  
    В безопасном распределении и хранении общего секретного ключа между участниками.

    \item \textbf{Как вычисляется функция Эйлера \(\varphi(n)\) для модуля \(n=pq\)?}  
    Для различных простых \(p\) и \(q\):
    \[
    \varphi(n)=\varphi(pq)=(p-1)(q-1).
    \]

    \item \textbf{Какие требования предъявляются к числу \(e\) при генерации открытого ключа?}  
    Должно выполняться \(1<e<\varphi(n)\) и \(\gcd(e,\varphi(n))=1\). Часто выбирают небольшое нечётное \(e\), например \(e=7\) или \(e=65537\).

    \item \textbf{Как найти число \(d\) для закрытого ключа, зная \(e\) и \(\varphi(n)\)?}  
    \(d\) находится как обратный элемент к \(e\) по модулю \(\varphi(n)\):
    \[
    d=e^{-1}\bmod\varphi(n),
    \]
    например, с помощью расширенного алгоритма Евклида.

    \item \textbf{Почему безопасность RSA зависит от сложности факторизации числа \(n\)?}  
    Если противник факторизует \(n=pq\), он вычислит \(\varphi(n)\), затем \(d\) и сможет расшифровывать сообщения. Эффективного алгоритма факторизации больших \(n\) в настоящее время нет.

    \item \textbf{Какие математические свойства обеспечивают корректность дешифрования в RSA?}  
    Соотношение \(ed\equiv1\pmod{\varphi(n)}\) и теорема Эйлера:
    \[
    m^{ed}\equiv m\pmod n.
    \]
    Для всех \(0\le m<n\) корректность также обеспечивается китайской теоремой об остатках.

    \item \textbf{Что произойдёт, если выбрать \(e\), не взаимно простое с \(\varphi(n)\)?}  
    Не существует обратного элемента \(d=e^{-1}\bmod\varphi(n)\), поэтому корректно сгенерировать секретный ключ и однозначно расшифровать сообщение нельзя.

    \item \textbf{Почему при выборе \(p\) и \(q\) из таблицы важно, чтобы они были простыми?}  
    Тогда \(n=pq\) имеет ровно два простых делителя, а функция Эйлера вычисляется как \((p-1)(q-1)\). Иначе формула \(\varphi(n)\) и корректность RSA нарушаются.

    \item \textbf{Объясните роль теоремы Эйлера в алгоритме RSA.}  
    Теорема Эйлера гарантирует, что при \(ed\equiv1\pmod{\varphi(n)}\) выполняется \(m^{ed}\equiv m\pmod n\). Это математически обосновывает восстановление исходного сообщения при дешифровании.
\end{enumerate}

\section*{Примечание}
Текст программы в отчёт не включён. Python-код находится в отдельном файле \texttt{lab2\_rsa.py}.

\end{document}